\subsection{How layers combine in each case}
\label{subsec:results-composition}

Combining the two layers in Porto Alegre, as in \texttt{flood+landslide}, does not add their effects. The number of pairs that lose access is the same as under \texttt{landslide} alone, 9{,}465, which confirms that the infinite factor dominates any finite one on a shared edge. What changes is the set of cells that remain connected and the cost of their paths. Under \texttt{landslide} the 2{,}446 cells that keep every destination pay nothing, because their shortest routes never cross an interdicted polygon. Restricting both scenarios to those same 2{,}446 cells, the median increase goes from 5.2 to 6.4 minutes and the 90th percentile from 49 to 81. The delay each cell already carried is compounded by the detour around the blocked polygons, and the 619 cells that lose access to every point of interest are counted separately, entering neither statistic.

The composite arrangement in Curitiba shows the same property from the other side. Chaining two simple rules would apply the vegetation discount to any tree-lined segment, including the ones where the thermal factor is already at its floor of 1, and to segments whose temperature was sampled at a different position along the edge. Reading both layers at identical points keeps the discount tied to the shade actually present where the heat was measured. The result is a reduction of the F15 loss from 2.4 to 1.8 points, measured against the typical scenario.
